\documentclass[11pt,a4paper]{article}
\usepackage[margin=1in]{geometry}
\usepackage[T1]{fontenc}
\usepackage[utf8]{inputenc}
\usepackage{lmodern}
\usepackage{microtype}
\usepackage{enumitem}
\usepackage{xcolor}
\usepackage{hyperref}
\hypersetup{colorlinks=true, linkcolor=black, urlcolor=blue}
\setlength{\parindent}{0pt}
\setlength{\parskip}{0.55em}
\setlist[itemize]{leftmargin=1.5em,itemsep=0.2em,topsep=0.2em}
\setlist[enumerate]{leftmargin=1.7em,itemsep=0.2em,topsep=0.2em}
\pagenumbering{arabic}
\begin{document}
\begin{center}
{\LARGE\bfseries Proposta para CRM Personalizado para Seguros e Planos de Saúde}
\end{center}
\vspace{0.5em}

Olá,\par

Meu nome é Fabricio Santana. Sou engenheiro de software e líder técnico com mais de 20 anos de experiência em desenvolvimento backend, arquitetura de software, integração com bancos de dados, testes automatizados, computação em nuvem e entrega de sistemas críticos. Também liderei equipes grandes e atuei em projetos em que consistência de dados, regras de negócio, segurança, validação e manutenção são fatores centrais.\par

Seu projeto tem forte aderência ao meu perfil porque o desafio principal não é apenas construir telas de cadastro. A parte importante é criar um CRM operacional confiável para seguros e planos de saúde, em que clientes, documentos, produtos, vendas, comissões, despesas, busca e indicadores funcionem juntos de forma segura, intuitiva e preparada para evolução.\par

Minha abordagem para esse tipo de projeto combina arquitetura web, modelagem de dados, controle de acesso, armazenamento de anexos, regras comerciais, dashboards, busca, responsividade, testes, implantação e handoff técnico em um plano de entrega coerente. Isso é importante aqui porque a qualidade do CRM depende de definir bem o primeiro release operacional e as regras de negócio antes de acelerar a implementação.\par

\section*{Metodologia de trabalho}

Meu método de entrega foi desenhado para reduzir risco antes do desenvolvimento principal. O trabalho começa com uma Discovery Sprint fechada e depois segue para um pacote de entrega somente quando escopo, prioridades, dependências, responsabilidades e critérios de aceite estiverem claros.\par

O primeiro passo recomendado é uma Discovery Sprint de 1 semana, com valor fixo de R\$ 3.000. Nessa fase, posso conduzir até 5 reuniões remotas para entender o fluxo comercial, revisar a gestão atual de clientes e vendas, mapear dados, produtos, comissões, anexos, indicadores, perfis de acesso, riscos, prioridades e o limite do primeiro release operacional. O entregável é um plano prático de implementação com escopo recomendado, abordagem técnica, marcos, critérios de sucesso, riscos, responsabilidades e recomendação do pacote de entrega.\par

Depois da discovery, o projeto pode seguir para um dos pacotes fixos de 4 semanas. O pacote é escolhido conforme complexidade, risco, número de módulos que entram no primeiro release, necessidade de segurança, regras de comissão, anexos, dashboard, busca, responsividade e implantação.\par

\textbf{Opções de pacote}\par

\begin{itemize}
\item \textbf{Essential Delivery} --- R\$ 12.000, 4 semanas. Inclui Scrum Master e 1 engenheiro full-stack. Adequado apenas se a discovery mostrar um recorte bastante estreito, com poucos cadastros, regras simples e baixa complexidade de dashboard.
\item \textbf{Product Acceleration} --- R\$ 24.000, 4 semanas. Inclui Scrum Master, 2 engenheiros full-stack e 1 engenheiro focado em testes, DevOps, observabilidade e deploy. Adequado para um primeiro release de CRM com clientes, vendas, produtos, anexos, dashboard, busca, responsividade, login e implantação com escopo controlado.
\item \textbf{Scale \& Intelligence} --- R\$ 36.000, 4 semanas. Inclui líder técnico, Scrum Master, 2 engenheiros full-stack e 1 engenheiro focado em testes, DevOps, observabilidade e deploy. Adequado quando o CRM exige permissões avançadas, auditoria, maior governança de dados, regras complexas de comissão, automações, integrações externas ou maior preocupação com segurança e escala.
\end{itemize}

Para esta oportunidade, a recomendação inicial mais provável é \textbf{Product Acceleration}, com possibilidade de revisão para \textbf{Scale \& Intelligence} se a discovery confirmar exigências mais fortes de LGPD, auditoria, múltiplos perfis, regras complexas de comissão, integrações ou automações.\par

As responsabilidades ficam claras desde o início: o Scrum Master organiza escopo, comunicação, checkpoints, acompanhamento de progresso e alinhamento de entrega; os engenheiros full-stack implementam backend, frontend, APIs, banco de dados, telas, regras de negócio e relatórios; o engenheiro de testes, DevOps, observabilidade e deploy apoia validação, CI/CD, configuração de ambiente, publicação, monitoramento e qualidade de release; e o líder técnico, quando incluído no pacote Scale \& Intelligence, responde por arquitetura, decisões técnicas, qualidade de código, segurança e gestão de risco.\par

\textbf{Garantia de defeitos de construção.} Após entrega e aceite, incluo garantia de 30 dias para defeitos de construção relacionados ao escopo aprovado. Isso cobre erros de implementação, comportamentos combinados que deixem de funcionar e defeitos introduzidos pelo código entregue. Não cobre novas funcionalidades, mudanças de regra comercial, novos campos ou relatórios não incluídos no escopo aprovado, alterações em serviços de terceiros, mudanças de infraestrutura fora do projeto ou solicitações que ampliem os critérios de aceite aprovados.\par

\section*{Análise da oportunidade}

Este projeto deve ser tratado como um CRM operacional para gestão comercial, e não apenas como um conjunto de telas. O sistema precisa ajudar a equipe a organizar clientes, acompanhar vendas, registrar produtos, controlar comissões, guardar documentos importantes e enxergar desempenho com indicadores claros.\par

Os módulos pedidos são coerentes, mas o risco real está em como eles se conectam:\par

\begin{itemize}
\item os dados de clientes precisam ter estrutura suficiente para o ramo de seguros e planos de saúde;
\item anexos precisam de organização, segurança, limite de tamanho, controle de acesso e estratégia de backup;
\item produtos, vendas, valores e comissões precisam seguir regras comerciais claras;
\item o dashboard precisa ter fórmulas bem definidas para vendas mensais, ganhos, despesas fixas e novos clientes;
\item a pesquisa precisa filtrar clientes e vendas de forma útil para a operação diária;
\item o uso mobile precisa priorizar os fluxos que realmente serão usados em campo ou em atendimento;
\item o sistema precisa considerar LGPD, principalmente se houver documentos contratuais, dados pessoais ou informações sensíveis associadas a planos de saúde.
\end{itemize}

Uma primeira versão prática deve focar em um release operacional controlado: cadastro de clientes, registro de produtos e vendas, anexos, cálculo inicial de comissões, dashboard essencial, busca e interface responsiva. Funcionalidades como automações, WhatsApp, funil avançado, importação complexa, integração com seguradoras ou operadoras e financeiro mais completo podem entrar depois, se forem priorizadas.\par

\section*{Plano de trabalho proposto}

Eu dividiria o trabalho nas seguintes fases:\par

\textbf{1. Discovery, fluxo comercial e desenho funcional}\par

\begin{itemize}
\item Mapear como a empresa hoje registra clientes, produtos, vendas, comissões, despesas e documentos.
\item Definir o primeiro release operacional e separar o que entra agora do que pode ficar para evolução.
\item Definir perfis de usuário, permissões, telas principais, dados obrigatórios e critérios de aceite.
\item Revisar cuidados de LGPD, anexos, backup, segurança e acesso mobile.
\item Produzir um plano de implementação com arquitetura, escopo, riscos, marcos e pacote recomendado.
\end{itemize}

\textbf{2. Arquitetura do CRM e base técnica}\par

\begin{itemize}
\item Definir a arquitetura web, autenticação, banco de dados, modelo de permissões e estratégia de deploy.
\item Estruturar entidades como clientes, contatos, produtos, vendas, comissões, despesas e anexos conforme o escopo aprovado.
\item Preparar a base para busca, filtros, relatórios e responsividade.
\item Definir tratamento de erros, validações, logs básicos e padrões para manutenção futura.
\end{itemize}

\textbf{3. Gestão de clientes, anexos, produtos e vendas}\par

\begin{itemize}
\item Implementar cadastro e edição de clientes com os campos definidos na discovery.
\item Implementar anexos vinculados a clientes ou vendas, com regras de organização e acesso.
\item Implementar cadastro de produtos e registro de vendas com valores, status e comissão associada.
\item Validar regras de cálculo e edição para evitar divergências comerciais nos relatórios.
\end{itemize}

\textbf{4. Dashboard, pesquisa e experiência mobile}\par

\begin{itemize}
\item Implementar dashboard com vendas mensais, ganhos totais, despesas fixas e novos clientes por mês.
\item Implementar filtros e pesquisa para clientes, vendas, produtos, status, período e responsáveis, conforme prioridade.
\item Ajustar os fluxos principais para uso em desktop, smartphones e tablets.
\item Validar leitura dos indicadores com dados reais ou amostras representativas.
\end{itemize}

\textbf{5. Testes, publicação, documentação e handoff}\par

\begin{itemize}
\item Testar cadastros, anexos, vendas, comissões, dashboard, filtros, permissões e responsividade.
\item Publicar a aplicação no ambiente acordado e orientar configuração inicial.
\item Documentar arquitetura, setup, regras de negócio, limites conhecidos e operação básica.
\item Fazer handoff técnico e orientação inicial para uso seguro do CRM.
\end{itemize}

\section*{Prazo estimado}

Minha recomendação imediata é 1 semana de Discovery Sprint. Depois disso, o primeiro ciclo de implementação pode seguir em 4 semanas, focado no primeiro release operacional acordado.\par

Cronograma prático sugerido:\par

\begin{itemize}
\item Semana 1: Discovery Sprint, mapeamento do fluxo comercial, modelo de dados, regras de comissão, indicadores, anexos, permissões, critérios de aceite e plano de implementação.
\item Semana 2: arquitetura, autenticação, banco de dados, base de clientes, produtos e vendas.
\item Semana 3: anexos, regras de comissão, validações, busca inicial e telas operacionais.
\item Semana 4: dashboard, filtros, responsividade e refinamento dos fluxos principais.
\item Semana 5: testes, ajustes, publicação, documentação, demo e handoff.
\end{itemize}

Se a discovery mostrar necessidade de importação complexa de dados, múltiplos perfis com auditoria, regras avançadas de comissão, integrações externas ou automações, a recomendação pode passar para Scale \& Intelligence ou para mais de um ciclo de implementação.\par

\section*{Disponibilidade e comunicação}

Posso atuar remotamente, com comunicação em português, por mensagens, chamadas agendadas e atualizações escritas. Prefiro checkpoints curtos ao longo da semana, especialmente após a definição do primeiro release, a validação do modelo de dados e os principais marcos de implementação.\par

Isso ajuda a manter o projeto alinhado com o fluxo real da operação comercial, as regras de comissão, os indicadores de gestão e os cuidados de segurança necessários para dados de clientes e documentos.\par

\section*{Orçamento}

Com base no entendimento atual, meu próximo passo recomendado é:\par

\begin{itemize}
\item \textbf{Discovery Sprint} --- R\$ 3.000, 1 semana.
\end{itemize}

Depois da discovery, o pacote mais provável para o primeiro release é:\par

\begin{itemize}
\item \textbf{Product Acceleration} --- R\$ 24.000, 4 semanas, caso o primeiro release inclua clientes, produtos, vendas, comissões, anexos, dashboard, busca, responsividade, login, permissões básicas e implantação.
\end{itemize}

Esta estimativa assume que o cliente fornecerá exemplos reais de dados, regras de venda e comissão, lista de produtos, definição de despesas fixas, exemplos de documentos, perfis de usuário, prioridade das telas e rapidez nas decisões sobre o primeiro release. Custos recorrentes de hospedagem, domínio, armazenamento de arquivos, serviços de terceiros, mensagens, e-mail ou integrações externas não estão incluídos no valor de desenvolvimento.\par

Se a revisão mostrar maior necessidade de LGPD, auditoria, múltiplos perfis, regras complexas, importação de dados, automações ou integração com ferramentas externas, recomendarei Scale \& Intelligence ou ciclos adicionais em vez de prometer um escopo maior do que cabe no primeiro release.\par

\section*{Perguntas antes da confirmação final}

Antes de confirmar escopo e marcos finais, eu gostaria de alinhar:\par

\begin{enumerate}
\item Quem usará o CRM no dia a dia: vendedores, gestores, administradores ou outros perfis?
\item Quais dados de cliente são obrigatórios no primeiro release?
\item O sistema precisa controlar apenas clientes e vendas ou também propostas, contratos, apólices, planos, renovações e dependentes?
\item Como as comissões são calculadas hoje: por produto, vendedor, venda, parcela, recorrência ou outra regra?
\item Quais despesas fixas devem aparecer no dashboard e como elas serão alimentadas?
\item Os anexos terão contratos, documentos pessoais, documentos assinados ou informações sensíveis?
\item Quais filtros de busca são indispensáveis para a operação diária?
\item Existem dados em planilhas, outro CRM ou sistema atual que precisam ser migrados?
\item O primeiro release precisa ter integração com WhatsApp, e-mail, seguradoras, operadoras, pagamentos ou outra ferramenta?
\item Já existe orçamento aprovado e data-alvo para a primeira versão em uso?
\end{enumerate}

Com essas respostas, consigo fechar o plano de entrega e a recomendação final de implementação.\par

\end{document}
